\documentclass{article}

\usepackage{amsmath,amssymb}
\usepackage{xurl}
\usepackage[hidelinks]{hyperref}
\setlength{\emergencystretch}{2em}

\input{command}

\title{Scope of Wolf Chapter~5 Operator Convexity and Monotonicity Results}
\author{}
\date{2026-08-24}

\begin{document}
\maketitle
\gapnote{scope-restriction}{open}

\section{The source assertions}

Wolf's Section~5.4 develops general results about operator-convex and
operator-monotone functions \cite[Section~5.4]{Wolf2012Quantum}. All operators
below act on finite-dimensional Hilbert spaces. The relation $A\leq B$ denotes
the Loewner order. A linear map $T$ is positive if it preserves positive
semidefiniteness, subunital if $T(\Id)\leq\Id$, and unital if
$T(\Id)=\Id$.

Wolf's Theorem~5.8 states that a function
$f\colon[0,\infty)\to\mathbb R$ is operator convex if and only if it has the
representation
\begin{align}
    f(x)
        &=a+bx+cx^2
          +\int_0^\infty\frac{yx^2}{y+x}\,d\mu(y),
    \label{eq:wolf_ch5_operator_convexity_scope_operator_convex_representation}
\end{align}
where $a,b\in\mathbb R$, $c\geq 0$, and $\mu$ is a finite positive measure
\cite[Theorem~5.8, Equation~(5.18)]{Wolf2012Quantum}.

L\"owner's theorem, stated as Wolf's Theorem~5.9, characterizes continuous
operator-monotone functions on $(0,\infty)$ by both their Pick-function
continuation and the representation
\begin{align}
    f(x)
        &=a+bx
          +\int_{-\infty}^0\frac{1+xy}{y-x}\,d\mu(y),
    \label{eq:wolf_ch5_operator_convexity_scope_operator_monotone_representation}
\end{align}
where $a\in\mathbb R$, $b\geq 0$, and $\mu$ is a finite positive measure
\cite[Theorem~5.9, Equation~(5.19)]{Wolf2012Quantum}.

Wolf's Theorem~5.10 assumes, in every dimension, that a Hermitian projection
$P$ and a Hermitian operator $A$ satisfy the compression inequality
\begin{align}
    f(PAP)&\leq Pf(A)P,
    \label{eq:wolf_ch5_operator_convexity_scope_source_projection_assumption}
\end{align}
whenever the spectrum of $A$ lies in an interval containing zero
\cite[Theorem~5.10, Equation~(5.22)]{Wolf2012Quantum}. It deduces
$f(0)\leq 0$, operator convexity of $f$, and the positive-subunital Jensen
inequality
\begin{align}
    f(T(A))&\leq T(f(A)).
    \label{eq:wolf_ch5_operator_convexity_scope_source_subunital_jensen}
\end{align}
The last formula is Equation~(5.23) of
Ref.~\cite{Wolf2012Quantum}.

Conversely, Wolf's Theorem~5.11 proves for an operator-convex function that
\begin{align}
    f(PAP)
        &\leq Pf(A)P+P^\perp f(0)P^\perp.
    \label{eq:wolf_ch5_operator_convexity_scope_source_projection_converse}
\end{align}
This is Equation~(5.25) of Ref.~\cite{Wolf2012Quantum}. When $f(0)\leq 0$, the
same theorem obtains
\begin{align}
    f(T(A))&\leq T(f(A))
    \label{eq:wolf_ch5_operator_convexity_scope_source_convex_jensen}
\end{align}
for positive subunital $T$; this is Equation~(5.26) of
Ref.~\cite{Wolf2012Quantum}. Wolf's Theorem~5.12 removes the condition
$f(0)\leq 0$ when $T$ is unital
\cite[Theorem~5.12]{Wolf2012Quantum}. Its entropic corollary is
\begin{align}
    T(\rho)\log T(\rho)
        &\leq T(\rho\log\rho).
    \label{eq:wolf_ch5_operator_convexity_scope_source_entropy_inequality}
\end{align}
This is Equation~(5.33) of Ref.~\cite{Wolf2012Quantum}.

For an operator-monotone function on $[0,a]$ satisfying $f(0)\geq 0$, Wolf's
Theorem~5.13 gives the reverse Jensen inequality
\begin{align}
    T(f(A))&\leq f(T(A)).
    \label{eq:wolf_ch5_operator_convexity_scope_source_monotone_jensen}
\end{align}
This is Equation~(5.34) of Ref.~\cite{Wolf2012Quantum}.

The corollary following Corollary~5.2 extends both Jensen inequalities to a
finite family of positive maps $T_1,\ldots,T_n$ satisfying
\begin{align}
    \sum_{i=1}^nT_i(\Id)&\leq\Id.
    \label{eq:wolf_ch5_operator_convexity_scope_finite_family_subunitality}
\end{align}
For operator-monotone $f$, it gives
\begin{align}
    \sum_{i=1}^nT_i(f(A_i))
        &\leq
        f\!\left(\sum_{i=1}^nT_i(A_i)\right).
    \label{eq:wolf_ch5_operator_convexity_scope_finite_family_monotone_jensen}
\end{align}
This is Equation~(5.37) of Ref.~\cite{Wolf2012Quantum}; the operator-convex
case gives the reverse inequality.

\section{The special-function restriction}

The formal development proves Jensen inequalities only for the functions
$x\mapsto x^p$ and $x\mapsto\log x$. For a positive subunital map $T$ and a
positive semidefinite operator $A$, it proves
\begin{align}
    T(A^p)&\leq(T(A))^p,
    &
    0&\leq p\leq 1,
    \label{eq:wolf_ch5_operator_convexity_scope_power_concavity}\\
    (T(A))^p&\leq T(A^p),
    &
    1&\leq p\leq 2.
    \label{eq:wolf_ch5_operator_convexity_scope_power_convexity}
\end{align}
For a positive unital map $T$ and a positive-definite operator $A$, it also
proves
\begin{align}
    T(\log A)&\leq\log(T(A)).
    \label{eq:wolf_ch5_operator_convexity_scope_logarithm_concavity}
\end{align}
The unital hypothesis in
Eq.~\ref{eq:wolf_ch5_operator_convexity_scope_logarithm_concavity} corrects
the false common subunital hypothesis in Wolf's Corollary~5.2
\cite[Corollary~5.2]{Wolf2012Quantum}.\footnote{The counterexample and
corrected statement are recorded in
\path{docs/paper-gaps/wolf_ch5_operator_jensen_lieb.tex}.}

The formal proofs use integral representations of the specified power
functions and derive the logarithmic case by a limit as $p\to 0^+$.\footnote{
Formal declarations:
\leanid{posMap_rpow_concave_jensen},
\leanid{posMap_rpow_convex_jensen}, and
\leanid{posMap_log_concave_jensen} in
\doclink{QICLean/Analysis/LiebConcavity}.}

Equations~\ref{eq:wolf_ch5_operator_convexity_scope_power_concavity},
\ref{eq:wolf_ch5_operator_convexity_scope_power_convexity},
and~\ref{eq:wolf_ch5_operator_convexity_scope_logarithm_concavity} are
instances of Wolf's general Jensen inequalities. They do not imply the
integral characterization of every operator-convex function, L\"owner's
theorem, the projection equivalences, the arbitrary-function Jensen
inequalities, or the finite-family corollary.

\section{Verdict}

The proved power and logarithm inequalities form a proper special-function
restriction of the results in Section~5.4 of Ref.~\cite{Wolf2012Quantum}. The
general operator-convexity and operator-monotonicity assertions remain an open
scope gap. No theorem for an arbitrary function $f$ follows from the three
special cases alone. The gap is resolved only by proving the general source
assertions themselves or source-faithful theorems that imply them.

\bibliographystyle{alpha}
\bibliography{references}

\end{document}
